%% LyX 2.4.4 created this file.  For more info, see https://www.lyx.org/.
%% Do not edit unless you really know what you are doing.
\documentclass[oneside,english]{amsart}
\usepackage[T1]{fontenc}
\usepackage[latin9]{inputenc}
\synctex=-1
\usepackage{mathrsfs}
\usepackage{enumitem}
\usepackage{amsthm}
\usepackage{amssymb}
\PassOptionsToPackage{normalem}{ulem}
\usepackage{ulem}

\makeatletter
%%%%%%%%%%%%%%%%%%%%%%%%%%%%%% Textclass specific LaTeX commands.
\newlength{\lyxlabelwidth}      % auxiliary length 
\theoremstyle{plain}
\newtheorem{thm}{\protect\theoremname}
\theoremstyle{plain}
\newtheorem{prop}[thm]{\protect\propositionname}
\theoremstyle{definition}
\newtheorem*{defn*}{\protect\definitionname}

%%%%%%%%%%%%%%%%%%%%%%%%%%%%%% User specified LaTeX commands.
\date{}
\usepackage{commath}

\makeatother

\usepackage{babel}
\providecommand{\definitionname}{Definition}
\providecommand{\propositionname}{Proposition}
\providecommand{\theoremname}{Theorem}

\begin{document}
\title{March~3 proposed changes: reordering sections}

\maketitle
These are overall suggestions I decided to type out. This is a proposed
outline of the entire paper. The main point of what's below is to
reorder what's already written. Anything appearing in double brackets
{[}{[}like this{]}{]} is an aside to you separate from what I'm proposing
for the paper.

\section{Mosaics.}

This section introduces the terms \emph{mosaic}, \emph{polygon mosaic},
and \emph{messy polygon mosaic}, probably in that order. It will include
what's now in Section~1 Introduction and Section~2 Related Work,
with the following things delayed. For all of these, we define them
for $\left(m,n\right)$, we should explicitly say that $m$ and $n$
are positive integers.

I'm wondering if a \emph{polygon mosaic} should include mosaics that
are entirely comprised of blank tiles. I'm leaning in this direction
because as I read the current draft, that choice just seemed to ``feel''
more and more correct. That would, unfortunately, make the term \emph{messy
polygon mosaic} weird. How attached are you to the term ``messy''?
I know it's been your terminology forever. We could instead define
\emph{polygon-free mosaic} (i.e., the complement of our current \emph{messy
polygon mosaic}) and count the number of those. Of course, the number
of mosaics that have at least one polygon will be $7^{mn}$ minus
the number of polygon-free mosaics. That's a big shift! I've taken
a couple days to think about it and it's growing on me. We should
discuss it.

Here are things I suggest omitting from this section:
\begin{itemize}
\item Don't actually define the matrices $A_{k}$, $B_{k}$, $C_{k}$, $D_{k}$,
nor state Theorem~1.1. We can do these things when we're ready to
prove the result. Rather, we say in the text that our goal is, for
any given positive integers $m$ and $n$, to determine the number
of polygon mosaics, and also the number of polygon-free mosaics (as
defined above). We refer to those two theorems below (e.g., ``We
will show how to calculate the number of polygon mosaics in Theorem~{[}ref{]}
and the number of polygon-free mosaics in Theorem~{[}ref{]}.'')
\item Minimal historical references (if any) and no mention at all of knots.
This all can come at the end of our paper. I'm not suggesting deleting
anything, just moving it.
\end{itemize}

\section{Binary lattices.}

Define binary lattice. I think that we no longer need $\left(m,n\right)$
binary lattices with $m$ or $n$ zero. That is, we only need to define
binary lattices for $m$ and $n$ positive integers.

Define \emph{good lattices} (as in my other document, i.e., not containing
$\begin{array}{cc}
0 & 1\\
1 & 0
\end{array}$ or $\begin{array}{cc}
1 & 0\\
0 & 1
\end{array}$ cells) and \emph{framed lattices} (aka framed binary lattices).

As in my other document, define the function $\mathscr{L}$, a bijection
from polygon mosaics to good framed lattices.

\section{Counting polygon mosaics.}

Here we count the number of $\left(m,n\right)$ polygon mosaics for
any positive integers $m$ and $n$. This might be a short section
(maybe it's a final part of the previous section??) involving only
defining $v$ (i.e., $0$ for $\begin{array}{cc}
0 & 1\\
1 & 0
\end{array}$ or $\begin{array}{cc}
1 & 0\\
0 & 1
\end{array}$, $1$ otherwise; then expanding to an arbitrary lattice) and saying
our desired number is $\sum v\left(\ell\right)$, where the sum is
over all good framed lattices.

\section{Matrix calculations.}

Motivated by the previous section we work with a more general $v$
and introduce matrices.

At this stage, $v$ is an arbitrary function from cells to the set
$\mathbb{C}$ of complex numbers. Then we extend $v$ to be a function
on binary lattices. We then define $A_{k}$, $B_{k}$, $C_{k}$, and
$D_{k}$ based on this arbitrary fixed $v$. 
\begin{prop}
For any $v$ as given above, $\sum v\left(\ell\right)$, where the
sum is over all $\left(m,n\right)$ good framed lattices, is equal
to the $\left(0,0\right)$ entry of $A_{n}^{m}$.
\end{prop}

{[}{[}Note: We have two big things to prove. One is above, and the
other below. Some of my motivation for the shuffling of the order
of topics is so that we can get a concrete result, i.e., for polygon
mosaics, without having to do both big proofs. Also, I think it's
nice to start with polygon mosaics, since polygon-free mosaics are
more complicated.{]}{]}

\section{Polygon-free mosaics.}

This will be where the inclusion exclusion principle is used (I suggest
an explicit statement, and my preference is to include a proof). Because
of the potential of changing from ``messy mosaics'' to ``polygon-free
mosaics,'' it's not clear which version of the inclusion exclusion
principle we want (i.e., $\left(-1\right)^{\left|J\right|}$ \uline{or}
$\left(-1\right)^{\left|J\right|-1}$, where $J$ runs over all subsets
of indices \uline{or} all nonempty subsets of indices).

Here we will define $v$ using the integers $0$, $1$, $-1$, and
$7$ in the usual way for messy mosaics (or polygon-free mosaics),
i.e., now $v$ does not just use $0$ and $1$ as above).

Here's a potentially useful rephrasing of stuff we already have: consider
when we have the size $\left(m,n\right)$ fixed and we enumerate all
polygons (my previous ``gridded polygons'') $M_{1}$, $M_{2}$,
\ldots , $M_{N}$. We consider the intersection $\bigcap_{j\in J}M_{j}$,
as $J$ varies through nonempty subsets of $\left\{ 1,2,\dots,N\right\} $.
Those $J$ such that $\bigcap_{j\in J}M_{j}$ is nonempty (i.e., the
$M_{j}$, $j\in J$ are ``compatible'') correspond one-to-one with
polygon mosaics. I think that's a useful intermediate step that we
don't currently have. Then those polygon mosaics correspond one-to-one
(via $\mathscr{L}$) with good framed lattices.

The intermediate step of talking about polynomial mosaics might be
nice because we can talk about all the mosaics that ``include''
a certain polynomial mosaic. More precisely, here's something we might
try.
\begin{defn*}
{[}{[}formerly ``compatible'' perhaps{]}{]} A mosaic $M$ \emph{includes}
the polynomial mosaic $P$ if $M$ can be formed from $P$ by placing
any tile (including the empty tile $T_{0}$) wherever $P$ has an
empty tile $T_{0}$.
\end{defn*}
\begin{prop}
Given a polynomial mosaic $P$, $\left|v\left(\mathscr{L}\left(P\right)\right)\right|$
is equal to the number of mosaics that include $P$, and the sign
of $v\left(\mathscr{L}\left(P\right)\right)$ is equal to $-1$ if
$P$ contains an odd number of polynomials, and $1$ if $P$ contains
an even number of polynomials.

{[}{[}Equivalent, not sure which is better{]}{]} For every good framed
lattice $\ell$, $\left|v\left(\ell\right)\right|$ equals the number
of mosaics that include $\mathscr{L}^{-1}\left(\ell\right)$ and the
sign of $v\left(\ell\right)$ equals $-1$ raised to the power of
the number of polygons in $\mathscr{L}^{-1}\left(\ell\right)$.
\end{prop}


\section{Extensions and reflections on earlier work}

Begin perhaps with extensions you and I know are simple but which
might not occur to readers. For example, if there are three types
of tiles connecting the midpoint of the top and bottom midpoints (e.g.,
the vertical dotted bar we already have plus one with a rightward
zig zag and one with a leftward zag zig). If we set $v\left(\begin{array}{cc}
0 & 1\\
0 & 1
\end{array}\right)$ and $v\left(\begin{array}{cc}
1 & 0\\
1 & 0
\end{array}\right)$ to $3$ then the techniques we used can be used to count the number
of polygon mosaics and the number of polygon-free mosaics. (Here the
number of ``messy mosaics'' would of course be $9^{mn}$ minus the
number of polygon-free mosaics.)

Then talk about work others have done, providing bounds on what we
found exactly (the number of polygon mosaics). The work on knots.
How our work can be easily adapted to replicate their results on the
number of knot mosaics. Conclude with an open question: how can we
calculate the number of knot-free mosaics. Of course, you and I struggled
with this for a long time but ultimately found that binary lattices
would probably not do the job. We can mention that we do not think
our techniques generalize to this open question.
\end{document}
